\newpage\pagestyle{introduction}

\unnumberedchapter{Introduzione}

\unnumberedsection{Contesto e motivazioni}

La qualità dei dati è un requisito operativo prima che teorico: un'analisi condotta su
informazioni incomplete produce conclusioni che nessun metodo statistico può recuperare. Quando
un attributo di una tupla è privo di valore, la scelta non è fra ``imputare'' e ``non imputare'',
ma fra imputare in modo informato e scartare il dato. È questa la ragione dell'interesse per la
\emph{data imputation}: ricostruire il valore mancante a partire dal resto dell'informazione
disponibile.

Tra gli approcci recenti, quelli basati su \emph{Relaxed Functional Dependencies} (RFD) hanno una
proprietà che li distingue: le regole di imputazione non vengono fornite dall'esterno, ma estratte
dal dataset stesso. Una RFD esprime un vincolo di dipendenza approssimato fra attributi ---
\emph{se due tuple concordano entro certe soglie sugli attributi di sinistra, allora concordano
entro una soglia sull'attributo di destra} --- e può quindi essere usata sia per generare
candidati all'imputazione sia per rifiutare quelli incoerenti. Il sistema \renuver{} \cite{breve2022renuver}
è un rappresentante di questa famiglia.

Il prezzo di questa ricchezza è computazionale: la fase di estrazione produce tipicamente migliaia
di regole, e il costo dell'imputazione cresce con il loro numero. Molte di esse, però, non
influenzano il risultato. Da qui il problema affrontato in questo lavoro: \emph{si possono
selezionare le RFD rilevanti senza alterare il comportamento del sistema?} È la domanda del
\emph{semantic pruning} applicato a un insieme di metadati.

\unnumberedsection{Problema affrontato}

Il sistema \renuver{} riceve in ingresso, oltre al dataset, un file di RFD. L'ipotesi di partenza
--- suggerita dalla struttura del sistema --- è che una parte di quel file sia ridondante e che
rimuoverla riduca il tempo di esecuzione senza degradare la qualità delle imputazioni.

Il problema si rivela però più sottile di così, per una ragione che questo lavoro porta alla luce
in modo formale: nel sistema studiato le RFD svolgono \emph{due ruoli distinti}, e i due ruoli
impongono requisiti opposti. Una RFD serve a \emph{generare} candidati all'imputazione, ma serve
anche a \emph{vetare} imputazioni che violerebbero il vincolo. Rimuovere una regola inutile per il
primo ruolo non garantisce che essa sia inutile per il secondo, e nulla nel formato del file
distingue i due casi.

Da questa osservazione discendono le due domande che strutturano la tesi:

\begin{enumerate}[label=\textbf{D\arabic*.},leftmargin=2.2em]
    \item \textbf{Che cosa si può potare con garanzia?} Esiste una classe di regole la cui
    rimozione è dimostrabilmente priva di conseguenze, e quanto è grande?
    \item \textbf{Che cosa si guadagna quando la garanzia si abbandona?} Se si rimuovono regole
    anche a costo di alterare il comportamento, quali sono i guadagni e a quali condizioni?
\end{enumerate}

\unnumberedsection{Domanda di ricerca}\label{sec:research-question}

Il lavoro muove da una domanda formulata in termini verificabili. Dato l'insieme di RFD che il
sistema riceve in ingresso, \emph{è possibile ridurne il numero almeno della metà mantenendo
l'accuratezza e la copertura entro cinque punti percentuali dalla versione originale, e riducendo
il tempo di esecuzione?} La domanda ha tre esiti distinti --- la riduzione, la qualità delle
imputazioni e il costo --- e la risposta del lavoro non è uniforme sui tre.

Sul primo esito la risposta è \emph{negativa}, e non per insufficienza delle strategie: il
\autoref{chap:theory} dimostra che una potatura che operi sul solo file delle regole e non rinunci
a nessuna garanzia non può superare, sui dataset esaminati, il \SI{20.9}{\percent} di riduzione.
Sul secondo esito la risposta è \emph{condizionata}: la potatura può migliorare il risultato --- su
Bridges lo migliora in modo statisticamente significativo --- ma soltanto quando la composizione
delle imputazioni lascia spazio al miglioramento, e sugli altri dataset danneggia. Sul terzo esito
la risposta è \emph{negativa nel caso generale e positiva in un caso}: la riduzione delle regole
non implica un risparmio di tempo, e la strategia che rimuove di più è quella che rallenta di più.
Il contributo del lavoro è di stabilire questi tre confini e di fornire il criterio con cui
riconoscere, su un dataset nuovo, in quale dei tre regimi ci si trovi.

La \autoref{tab:research-question} anticipa la forma delle tre risposte, con il rimando alle
sezioni in cui ciascuna è argomentata: serve a rendere esplicito fin dall'inizio che cosa il
lavoro sosterrà e che cosa non sosterrà.

\begin{table}[htbp]
\centering
\caption{Risposta alla domanda di ricerca, per esito. La tabella anticipa la struttura
dell'argomentazione; le grandezze sono discusse nelle sezioni indicate.}
\label{tab:research-question}
\footnotesize
\begin{tabular}{@{}p{2.5cm}p{7.3cm}p{2.4cm}@{}}
\toprule
esito & risposta & dove \\
\midrule
riduzione $\ge \SI{50}{\percent}$ \emph{con garanzia} & \textbf{impossibile}: la classe di regole
la cui rimozione è certificabile non supera il \SI{20.9}{\percent} sui dataset esaminati, e il
limite è del criterio, non dell'implementazione & \autoref{chap:theory}, \autoref{sec:maxextra} \\
accuratezza e copertura entro cinque punti & \textbf{condizionata}: il risultato migliora dove la
composizione delle imputazioni ha spazio per migliorare e la copertura non cala, e peggiora dove
una delle due condizioni manca & \autoref{sec:res-a5b}, \autoref{sec:res-boundary} \\
riduzione del tempo di esecuzione & \textbf{non garantita}: il tempo dipende dalla composizione
dell'insieme potato, non dalla sua dimensione; il caso migliore accelera di \num{11} volte, il
peggiore rallenta di \num{251} & \autoref{sec:res-cost} \\
\midrule
\multicolumn{3}{@{}p{12.2cm}@{}}{\footnotesize Il criterio diagnostico della
\autoref{sec:diagnostic} rende operativa la seconda risposta: decide a costo nullo se la leva del
miglioramento esiste, e impone di misurare la copertura prima di adottare la potatura.} \\
\bottomrule
\end{tabular}
\end{table}

\unnumberedsection{Obiettivi del lavoro}

Gli obiettivi, in ordine di priorità:

\begin{itemize}[leftmargin=1.4em]
    \item \textbf{Caratterizzare il sistema.} Derivare dal comportamento effettivo di \renuver{}
    --- non dalla sua documentazione --- le condizioni sotto le quali una RFD può influenzare
    un'imputazione, e tradurle in criteri di potatura verificabili.
    \item \textbf{Definire una metodologia di pruning.} Progettare e implementare strategie di
    selezione, distinguendo esplicitamente quelle che offrono una garanzia da quelle euristiche.
    \item \textbf{Valutare l'impatto.} Misurare l'effetto della potatura sulle prestazioni
    dell'imputazione --- accuratezza, copertura, precisione e tempi --- con un protocollo
    appaiato e test statistici, su più famiglie di dataset.
    \item \textbf{Spiegare i risultati.} Non limitarsi a misurare: individuare il meccanismo che
    determina quando la potatura conviene e quando danneggia.
\end{itemize}

\unnumberedsection{Contributi della tesi}

I contributi originali di questo lavoro sono i seguenti.

\begin{enumerate}[leftmargin=1.6em]
    \item \textbf{Tre condizioni esatte derivate dal bytecode} \Vone{}, \Vtwo{}, \Vthree{},
    che classificano le RFD in base al loro ruolo effettivo nel sistema. \Vtwo{} è la più
    rilevante: identifica le regole il cui potere di veto è implicato da un'altra regola
    sopravvissuta, e la loro rimozione è quindi priva di effetti.

    \item \textbf{Un teorema negativo} (\autoref{chap:theory}) che delimita il campo: se la
    potatura opera sul solo file delle regole e ogni rimozione è certificata con l'informazione
    ivi disponibile, allora le regole rimovibili sono quelle inerti oppure quelle il cui veto è
    implicato da una regola sopravvissuta. Il risultato è quantificato: la classe potabile con
    garanzia non supera il \SI{20.9}{\percent} sui dataset studiati, e si dimostra che tale tetto
    è \emph{intrinseco al criterio} e non un artefatto della sua implementazione
    (\autoref{sec:maxextra}). Il teorema separa ciò che è dimostrato --- la caratterizzazione
    delle rimozioni certificabili --- da ciò che è misurato, e la distinzione è mantenuta in
    tutto il lavoro.

    \item \textbf{Una verifica empirica della condizione \Vtwo{}} (\autoref{sec:res-vetocover})
    su \num{1242321} terne con zero
    discrepanze, e la conferma che la sua applicazione lascia l'output identico al sistema
    originale su \num{113} configurazioni su \num{113}, su quattro famiglie di dataset.

    \item \textbf{Un'euristica di potatura aggressiva} (\str{A5b}, \autoref{sec:res-a5b}) che su Bridges rimuove
    l'\SI{84.2}{\percent} delle regole \emph{migliorando} il risultato (\num{69} imputazioni
    corrette in più, \num{72} errori in meno, \num{11} volte più veloce), e la misurazione del
    fatto che la stessa strategia danneggia gli altri dataset.

    \item \textbf{La spiegazione del perché} (\autoref{sec:mechanism}). Il guadagno della potatura aggressiva si scompone
    nel prodotto di due fattori --- la quota di imputazioni ottenute per via esatta e il numero
    di imputazioni effettuate. La potatura conviene solo quando entrambi crescono. Ne discende
    un criterio diagnostico misurabile prima di potare, che sostituisce la soglia euristica
    precedentemente adottata.

    \item \textbf{Sei difetti documentati} (\autoref{chap:defects}), quattro del materiale
    d'archivio e due del codice sviluppato, con impatto quantificato e procedura di
    riproduzione: con procedura di
    riproduzione: un ramo irraggiungibile nel codice di calcolo della similarità, un
    disallineamento di schema nei dataset Glass, indici di ground truth non ricostruibili nei
    dataset Physician, e una build con istruzioni di debug. A questi si aggiunge un difetto di
    correttezza nel codice sviluppato in questo lavoro, individuato e corretto.

    \item \textbf{La caratterizzazione del costo computazionale} del sistema
    (\autoref{sec:cost}), ottenuta ricompilando il codice originale da sorgenti decompilati ---
    con equivalenza verificata rispetto all'eseguibile di partenza --- e strumentandolo: il
    \SI{99}{\percent} del tempo di esecuzione è nella verifica dei veti, non nella generazione
    dei candidati.

    \item \textbf{La confutazione della proporzionalità fra riduzione e tempo}
    (\autoref{sec:res-cost}). Il lavoro mostra che il tempo di esecuzione dipende dalla
    \emph{composizione} dell'insieme potato e non dalla sua dimensione: la strategia che rimuove
    fino al \SI{98}{\percent} delle regole rallenta l'esecuzione fino a \num{251} volte, e una
    variante che conserva più regole di un'altra risulta circa nove volte più lenta. La
    spiegazione candidata --- la perdita della scorciatoia esatta per le tuple con una sola cella
    nulla --- è dichiarata come non verificata.
\end{enumerate}

\unnumberedsection{Struttura dell'elaborato}

Il \autoref{chap:background} introduce i concetti preliminari e colloca il lavoro rispetto alla
letteratura: qualità dei dati, valori mancanti, tecniche di imputazione, \emph{Relaxed Functional
Dependencies}, riduzione di un insieme di regole, misure di similarità e la nozione di semantic
pruning adottata.

Il \autoref{chap:renuver} analizza il sistema \renuver{} per disassemblaggio: architettura, flusso
di imputazione, dataset impiegati, il doppio ruolo delle RFD, le due procedure di riempimento e la
caratterizzazione del costo computazionale.

Il \autoref{chap:methodology} presenta la metodologia: motivazioni, spazio delle strategie, le
strategie implementate, le condizioni esatte di potatura derivate dal comportamento del sistema e
il contratto di integrazione.

Il \autoref{chap:theory} formula e dimostra il risultato negativo che delimita il campo di
applicabilità, ne trae la conseguenza sul target di riduzione e ne precisa la portata.

Il \autoref{chap:hypotheses} espone come sono nate le ipotesi, le sette ipotesi primarie con il
loro criterio di falsificazione e le domande esplorative.

Il \autoref{chap:experiments} descrive la valutazione sperimentale: il protocollo di verifica e
la metodologia statistica, l'infrastruttura, i dataset e le configurazioni, le metriche, il
termine di paragone banale, gli esperimenti condotti con il collegamento alle ipotesi e lo studio
di ablazione.

Il \autoref{chap:results} analizza i risultati: la potatura che non perde nulla, il confronto
con la riduzione di ridondanza della teoria classica, la potatura che produce il guadagno e il
confine in cui il guadagno non si realizza, il meccanismo che li spiega, il criterio diagnostico,
i risultati negativi e il costo.

Il \autoref{chap:discussion} discute gli esiti delle ipotesi, il quadro complessivo, i limiti
dell'evidenza e la collocazione del lavoro rispetto alla letteratura.

Il \autoref{chap:defects} raccoglie i difetti e le criticità: quelli del sistema studiato, quelli
del codice sviluppato in questo lavoro e le conseguenze generali che se ne traggono.

Il \autoref{chap:conclusions} sintetizza il lavoro e ne indica gli sviluppi possibili. L'\autoref{app:reproducibility} documenta la riproducibilità dei risultati e
l'\autoref{app:results} raccoglie le tabelle complete delle campagne.
